\documentclass[UTF8,11pt]{ctexart}

\usepackage[margin=1in]{geometry}
\usepackage{microtype}
\usepackage{amsmath,amssymb,amsthm,mathtools}
\usepackage{booktabs,tabularx,array,multirow}
\usepackage{graphicx}
\usepackage{caption}
\usepackage{subcaption}
\usepackage{placeins}
\usepackage{xcolor}
\usepackage{enumitem}
\usepackage{natbib}
\usepackage[hidelinks]{hyperref}
\usepackage{fancyhdr}
\usepackage{seqsplit}

\definecolor{navy}{HTML}{17365D}
\definecolor{teal}{HTML}{00897B}
\hypersetup{colorlinks=true,linkcolor=navy,citecolor=teal,urlcolor=navy}
\graphicspath{{figures/}}
\captionsetup{font=small,labelfont=bf}
\setlist{nosep,leftmargin=1.8em}
\setlength{\parskip}{0.35em}
\setlength{\parindent}{2em}
\pagestyle{fancy}
\fancyhf{}
\fancyhead[L]{前缀因果损害预算约束的跨域电池健康状态估计}
\fancyhead[R]{吴雨阳}
\fancyfoot[C]{\thepage}
\renewcommand{\headrulewidth}{0.4pt}
\setlength{\headheight}{14pt}

\newtheorem{lemma}{引理}
\newtheorem{theorem}{定理}
\newtheorem{corollary}{推论}
\newcommand{\R}{\mathbb{R}}

\title{\textbf{基于前缀因果损害预算投影的\\风险可控跨域锂离子电池健康状态估计}}
\author{吴雨阳}
\date{}

\begin{document}
\maketitle

\begin{abstract}
跨域电池健康状态（state of health，SOH）模型通常必须在新车队尚无容量标签时给出预测。此时，源域学习更新可能在未知域产生负迁移，回顾性单调平滑还可能在未来测量到达后修改先前估计。本文提出前缀因果损害预算投影（prefix-causal harm-budget projection，PCHP），在输出空间中分离受保护源域预测器与学习候选。PCHP 逐记录构造单向因果基线状态，并将候选投影到声明损害管、物理范围与递归可行非增集合的交集。对任意可能真实值，相对于受保护状态的绝对损失增量均不超过预算 $\delta$，且无须分布或独立性假设。绝对损失管是精确可行域；零预算禁止任何结果未知的非平凡更新，而 PCHP 是完整在线约束下唯一最小化候选失真的解。损失几何扩展进一步表明：度量损失对应精确度量球；有界区间上的平方损失具有依赖基线位置的精确损害域，而当结果遍历整个实数轴时，所有非平凡有限预算更新均不可能。严格嵌套留一完整域评估覆盖 $12$ 个公开电池域、$586$ 个电芯和 $601{,}932$ 条记录。PCHP 将域等权电芯宏平均 MAE 从 $0.07524$ 降至 $0.07038$；配对差为 $-0.00486$，域聚类 $95\%$ 区间为 $[-0.00671,-0.00301]$，并在 $12$ 个域中的 $11$ 个改善。在 $\delta=0.01$ 时，PCHP 相对于源域选择的候选无关安全偏移进一步降低 MAE $0.00477$；协议锁定的回顾性预算路径得到预算归一化积分差 $-0.00694$，区间为 $[-0.01108,-0.00296]$。更强的源域调优无保护因果比较器达到 $0.05470$，却在所有域违反 $\delta=0.01$ 损害管，揭示安全—精度机会成本。结果隔离的拜罗伊特公开数据确认覆盖 $45$ 个合格电芯：电芯等权 MAE 从 $0.17496$ 降至 $0.16510$，配对差为 $-0.00986$，电芯 bootstrap 区间为 $[-0.01000,-0.00959]$，且 $45/45$ 个电芯均改善。由于该域存在系统性低估，事后固定上移 $0.01$ 的精度略高于 PCHP，因此外部结果验证的是冻结损害—效用契约，而非自适应优越性。另一项 NASA 压力测试保留了确定性证书，同时暴露绝对目标不兼容区域。因此，PCHP 提供的是可审计的风险有界更新契约，而非基线准确性或普遍部署安全性的证书。
\end{abstract}

\noindent\textbf{关键词：}锂离子电池；健康状态；跨域学习；前缀因果；确定性损害界；单调预测；外部确认

\section{引言}

设想一个正在引入新供应商电芯的换电运营商。投入运行后的最初数月，电池管理系统持续接收部分充电记录，但准确容量标签仍不可得，因为完整参考放电会中断车辆运营。源域模型能够立即估计 SOH，另一模型则可以利用电芯相对于投运初期的充电响应变化提出更新。然而，在新车队结果出现之前，更新的方向和幅度都无法验证。更严重的是，如果通过回顾性单调回归施加“SOH 长期不升”的结构约束，下个月出现的一次向下异常就可能修改今天已经发布的健康记录。即使最终重构轨迹十分平滑，这样的估计器也不能作为在线维护记录。

SOH 估计直接影响维护与质保决策。数据驱动模型能够利用电压、电流、充电量、温度和循环历史，而无须辨识每一条退化路径 \citep{ng2020mlbattery,roman2021pipeline}；短充电片段可用于 SOH 估计，早期寿命记录则可在完成全寿命表征之前支持寿命预测 \citep{roman2021pipeline,severson2019cyclelife}。跨域学习进一步通过表征对齐、源域选择、个性化迁移、模型集成和测试时适应降低对目标域老化标签的依赖 \citep{li2020sstca,ma2022personalized,lu2023without,duan2025wasserstein,feng2024batteryttt}。这些进展回答了如何产生有用更新，却没有自动回答两个部署问题：对于每一个隐藏真实值，更新相对于受保护预测可能造成多大损害；以及轨迹约束能否在不访问未来记录的条件下执行。

本文通过分离“效用产生”“因果状态构造”与“损害控制”来填补这一能力缺口。受保护模型使用原始早期充电特征，学习候选进一步使用相对于前五次表征记录的原量纲变化。受保护原始预测被转换为忽略向上创新、按比例 $\alpha$ 吸收向下创新的前缀因果非增状态；随后，PCHP 将当前候选投影到该状态周围的递归可行区间。整个过程不使用目标结果，不修改已经发布的前缀，并限制相对于受保护状态的最坏绝对损失增量。图~\ref{fig:workflow_zh} 概括了现实问题、方法与证据链。

\begin{figure}[t]
  \centering
  \includegraphics[width=\textwidth]{fig1_method_workflow_zh.pdf}
  \caption{现实问题、前缀因果方法与核心证据。\textbf{A}，在线 SOH 记录在未来充电数据到达后不得改变。\textbf{B}，受保护原始预测被转换为单向因果非增状态，投运参照模型提出候选更新。\textbf{C}，PCHP 将候选投影到损害管与递归可行非增集合的交集。\textbf{D}，嵌套开发评估、候选无关控制和协议冻结 NASA 压力测试的配对效应；区间按所述独立单位重采样。}
  \label{fig:workflow_zh}
\end{figure}

本文只有一个核心贡献：提出一个递归可行的输出空间投影，使前缀执行、非增输出和相对于受保护预测的逐记录损害预算同时成立。理论进一步证明，损害管是给定绝对损失预算下必要且充分的最大可行输出集合；若要求无条件零损害，则除保持受保护预测不变之外不存在非平凡更新；PCHP 还是完整在线约束下唯一最接近学习候选的解。理论还揭示损失依赖的可行性边界：度量损失产生精确度量球；有界 SOH 范围上的平方损失产生依赖基线位置的精确非对称区间；结果遍历整个实数轴时，平方损失则使所有非平凡有限预算更新均不可能。一个方向加权推论进一步把低估与高估的不同代价转换为精确非对称损害管；当代价和预算随记录变化时，同一分析还给出精确递归可行条件与紧的当前最小预算。这些结果仍属于同一个输出契约贡献。投运参照候选是该契约内部的效用机制，而不是另一个并列的架构创新。本文通过五层相互衔接的证据验证该贡献：未来扩展证伪证明 PCHP 从不修改历史；十二域严格嵌套评估及候选无关安全偏移控制；揭示预算机会成本的源域调优无保护比较器；预测先于容量访问冻结的拜罗伊特结果隔离确认；以及在放电协议导致绝对标签不可比时保留完整失败边界的 NASA 压力测试。

该保证是相对保证。PCHP 不保证受保护基线本身准确，不保证候选在每个数据域都改善，不将非增 SOH 视为无短时容量恢复的电化学定律，也不保证 $\delta=0.01$ 对所有运营决策均合适。它只证明：已经发布的更新具有前缀因果性、保持非增，并且相对于受保护因果状态的绝对损失增量不超过声明预算。

\section{相关工作与原创性边界}

\subsection{跨域电池健康状态估计}

电池 SOH 估计已经从人工健康指标和浅层回归器发展至序列模型、物理信息学习和跨域迁移 \citep{ng2020mlbattery,roman2021pipeline,wang2024pinn}。半监督迁移成分分析通过对齐源域与目标域表征来缓解有限标签问题 \citep{li2020sstca}；深度迁移学习面向未见放电协议开展个性化健康预测 \citep{ma2022personalized}；域适配模型集成则在不新增目标域退化实验的条件下估计 SOH \citep{lu2023without}。近期研究还使用 Wasserstein 距离选择源域后再微调时序模型 \citep{duan2025wasserstein}，或利用逐步到达的无标签目标输入执行测试时训练 \citep{feng2024batteryttt}。同期工作进一步探索使用少量目标监督的域对抗电芯到电池包迁移 \citep{shu2026celltopack}、稀疏目标标签下的概率迁移 \citep{shao2026probabilistic}，以及结合全局单调惩罚的物理引导专家路由 \citep{jiang2026phasemoe}。面向可靠性的扩展还将域对齐与共形预测区间结合 \citep{filgueira2026conformalized}，或根据带标签现场留出集上的误差选择单调后校准器 \citep{nyachionjeka2026safecalibrated}。这些方法改进了适应能力、不确定性量化或物理一致性，但并未给出本文研究的保证。本文的目标域既不提供结果标签，也不参与参数更新：两个源域模型在目标预测前已经冻结，目标阶段只约束其输出轨迹，保证对象是逐点损失增量，而不是区间覆盖率、目标监督改善或训练惩罚。

相邻的预测与健康管理研究使上述区别更加明确。复合域偏移预测同时对齐机械退化轨迹的域间分布和域内退化阶段分布 \citep{ding2023compound}；源域开放集域泛化在训练阶段无目标域数据的条件下，面向未知工况学习故障分类器 \citep{song2026regroupment}；交互式电池预测则递归耦合神经网络容量估计与切换状态空间过程以更新剩余使用寿命 \citep{duan2025interactive}。这些工作的报告目标分别是降低表征偏移、执行仅用源域的开放集分类以及递归寿命预测。PCHP 针对不同的保证对象：在学习预测器保持冻结的条件下，围绕受保护的前缀因果状态约束与结果无关的逐记录损害预算。

BatteryML 和 BatteryLife 对异构公开老化数据进行了标准化整理，从而使大范围跨数据集评估更易复现 \citep{zhang2024batteryml,tan2025batterylife}。但数据广度也要求正确界定统计重复：循环嵌套于电芯，电芯又嵌套于完整数据集。本文因此将完整电池数据域而非单个循环作为开发阶段跨域效应的独立单位。

\subsection{负迁移与安全回归}

负迁移指迁移数据或知识降低目标性能的情形，既有缓解路线包括域相似性估计、安全迁移和显式负迁移抑制 \citep{zhang2023negative}。在安全半监督回归中，SAFER 将监督预测几何投影至多个半监督回归器构成的凸集，并在完整真实标签向量位于该集合时获得平方损失不劣结论 \citep{li2017safer}。SAFEW 将极大极小安全弱监督扩展至多种分类与回归凸损失，但仍要求完整真实标签向量与基学习器凸包满足相应条件 \citep{li2021safeweak}。专家评价者建议允许不同专家及不同时刻采用不同的适当可混合损失，但其结论是在结果到达后成立的累积遗憾界 \citep{chernov2009evaluators}；在线概率再校准同样追求校准性以及相对于基线预测器渐近消失或较低的累积遗憾 \citep{marx2025calibrated,deshpande2025calibrated}。加权共形迁移利用无标签目标协变量，在协变量偏移模型下恢复边际覆盖率 \citep{tuwani2023transport}；电池领域的共形迁移则把表征对齐与校准预测区间结合 \citep{filgueira2026conformalized}。受保护在线熵匹配通过检测分类熵偏移并更新模型参数，使其在无分布偏移时经验上保持准确率与校准 \citep{bar2024protected}。带标签的电池现场数据也被用于仅在留出误差改善时选择单调后校准器 \citep{nyachionjeka2026safecalibrated}。这些都是实质性的保护机制，但其保护对象分别是条件式向量风险、累积遗憾、边际覆盖率、经验留出误差或自适应训练过程，而不是任意隐藏真实值下每条已发布点预测的损失。

本文的区别在于保证对象和执行契约。本文不声称最终预测必然不劣，而是在不要求真实值位于候选集合的条件下，对每一个实数真实值、每一条已发布记录给出绝对损失最大增量的精确上界。该保证在真实值被观察之前即可核验，两个预测器均保持冻结，而且同一投影必须保持此前已经发布的全部非增前缀。表~\ref{tab:novelty_audit_zh} 直接比较这些信息边界与保证边界。

\begin{table}[t]
  \centering
  \caption{按照保证对象与信息边界开展的最近邻方法审计。“条件”仅记录相应文献陈述保护结论时的主要条件，不评价方法优劣。}
  \label{tab:novelty_audit_zh}
  \small
  \begin{tabularx}{\textwidth}{@{}>{\raggedright\arraybackslash}p{0.17\textwidth}>{\raggedright\arraybackslash}p{0.20\textwidth}>{\raggedright\arraybackslash}p{0.25\textwidth}>{\raggedright\arraybackslash}X@{}}
    \toprule
    方法 & 保护对象 & 信息或成立条件 & 与 PCHP 的区别 \\
    \midrule
    SAFER / SAFEW \citep{li2017safer,li2021safeweak} & 相对监督基线的完整向量损失 & 真实标签向量位于基回归器凸包 & 离线向量或极大极小预测；没有逐记录结果一致预算或轨迹递归 \\
    专家建议 / 在线再校准 \citep{chernov2009evaluators,marx2025calibrated,deshpande2025calibrated} & 校准性与累积损失遗憾 & 结果依次到达并更新学习器或校准器 & 累积保证，而非结果观察前的逐点损害域 \\
    共形迁移 \citep{tuwani2023transport,filgueira2026conformalized} & 预测集合的边际覆盖率 & 无标签目标协变量及所声明的共形偏移假设 & 保护集合覆盖率，而非已发布点预测相对基线的损失 \\
    受保护适应 \citep{bar2024protected} & 无偏移数据流中的准确率与校准 & 目标熵数据流更新模型参数 & 围绕训练过程的经验保持，而非冻结模型的确定性逐点证书 \\
    电池安全校准 \citep{nyachionjeka2026safecalibrated} & 目标留出误差与单调尺度 & 带标签目标留出集选择校准器 & 目标监督后校准；没有结果未知条件下的逐点证书 \\
    PCHP（本文） & 相对受保护因果状态的逐记录绝对损失增量 & 不要求目标标签、结果分布或独立性 & 精确预算管与前缀因果、递归非增可行性联合成立 \\
    \bottomrule
  \end{tabularx}
\end{table}
\FloatBarrier

因此，本文的原创性不来自标量裁剪、几何投影或对负迁移现象的发现。本文经核验的区别是一个面向电池部署的信息结构组合：仅由源域学习的候选、前缀因果受保护状态、递归可行非增输出以及精确逐点绝对损失损害管。在本次检索边界内没有发现完全相同的组合，但这只能支持限定区别，不能证明普适意义上的“首次”。

损失几何扩展同样保持这一边界。度量恒等式来自既有的度量性质，本文不把它作为新的数学原语。更窄的区别是：给出平方损失关于结果范围的精确可行性分界，并将其接入同一前缀因果递归契约；安全弱监督方法在真值集合条件下保护向量损失，而专家建议与在线再校准在结果到达后保护累积损失。

递归可行性本身是约束控制中的既有概念，其强形式要求对每一条可行输入序列而非仅对优化器选中的序列均保持后续可行 \citep{lofberg2012oops}。预测控制与在线优化也已经研究时变约束和依赖历史的可行集合 \citep{li2021information,garber2024projectionfree,knaup2026recursively}。这些工作的对象是动态状态、累计约束违反、动态遗憾或机会约束，并不提供本文研究的未知结果一致逐点损失域。因此，本文不把递归可行性或时变约束本身作为原创性主张；本文更窄的增量是由一维 PCHP 损害管递归诱导出的精确条件。

\subsection{单调退化趋势}

长期容量衰减支持对非增趋势施加结构约束，但实测容量也可能由于静置恢复等因素出现局部回升和波动 \citep{ma2021regeneration}。在线等距回归逐次输出预测，并以事后最优单调函数为比较对象分析在线遗憾，还可扩展至平方损失之外 \citep{kotlowski2016online}；其目标是对抗式在线遗憾，而不是围绕电池特定基线提供损害保护。电池文献也常将基于短充电或放电片段的即时推断称为“在线 SOH 估计” \citep{liu2020online}，但特征在线可用并不自动保证后续轨迹平滑器不会修改历史输出。PCHP 处理的是这一更窄的信息边界。其非增状态用于估计潜在长期趋势，可能抑制真实的短时容量恢复，因此原始测量与受保护趋势均应被保留并可审计。

\section{问题定义}

设 $i$ 表示物理电芯，$t$ 表示其有序表征记录。容量型目标定义为
\begin{equation}
  y_{it}=\frac{Q_{it}}{\operatorname{median}(Q_{i1},\ldots,Q_{i5})},
  \label{eq:soh_zh}
\end{equation}
其中，前五次有效表征容量的中位数构成电芯特定的结果参照。在预测阶段，$Q_{it}$ 与 $y_{it}$ 均不可用。预测从前五次充电记录之后开始；这些记录的充电测量而非放电容量同时构成投运特征参照。一个数据域对应同一实验来源、仪器体系、化学体系或协议所生成的一组电芯。在每个开发折中，完整目标域均排除于模型拟合和超参数选择之外。

令 $\mathcal F_{it}$ 包含冻结的源域模型、电芯前五次充电参照，以及截至记录 $t$ 已观测到的全部充电测量。若某个在线预测关于 $\mathcal F_{it}$ 可测，且在 $t$ 之后追加任意记录都不会改变已经发布的前缀，则称其具有前缀因果性。记 $r_{it}$ 为受保护模型的原始预测，$c_{it}$ 为学习候选，两者均关于 $\mathcal F_{it}$ 可测。最终输出 $p_{it}$ 需要同时满足四项要求：前缀因果、轨迹非增、物理取值范围 $\mathcal Y=[0,1.3]$，以及相对于受保护因果状态 $b_{it}$ 的绝对损失损害预算 $\delta$。

受保护状态是一项参照契约，而不是最优性断言。它的准确性决定 PCHP 可以安全活动的邻域；下述定理不能把失准基线或协议不兼容的标签转化为准确估计。

\section{方法}

\subsection{早期充电特征与投运参照原量纲变化}

每条表征记录由充电开始后的前 $600\,\mathrm{s}$ 数据构造。$24$ 个原始通道包括归一化累计充电量、电压均值与标准差、电压斜率、电流统计量、温度可用性与温度统计量、每隔 $60\,\mathrm{s}$ 的电压，以及 $0$、$300$ 和 $600\,\mathrm{s}$ 时的电流；模型另使用对数变换后的循环序号。缺失值以源训练折中的中位数填补，并由填补器附加缺失指示变量。

对于第 $j$ 个原始通道，电芯 $i$ 的投运参照为
\begin{equation}
  a_{ij}=\operatorname{median}_{1\leq s\leq5}x_{ijs},
\end{equation}
原量纲有符号变化特征定义为
\begin{equation}
  d_{ijt}=x_{ijt}-a_{ij}.
  \label{eq:change_zh}
\end{equation}
受保护模型仅接收原始特征和循环序号；候选模型接收相同输入及全部 $d_{ijt}$。开发阶段曾评估归一化相对变化，但其域等权误差更高，且在参照量接近零时会放大分母噪声，因此在外部数据访问之前予以删除。该结论被明确报告为开发消融，而非独立外部发现。

两个主预测器均采用 ExtraTrees 回归，包含 $80$ 棵树，最小叶节点样本数为 $20$，特征子采样比例为 $0.8$，随机种子固定为 $20{,}260{,}801$。每个源域物理电芯最多保留 $80$ 个沿时间均匀抽取的记录。受保护基线采用普通合并样本权重；候选模型则使每个源数据域具有相同的总权重，避免大型档案仅因循环数量更多而主导模型。目标域推理过程中不使用结果标签，不进行微调或源域选择。

\subsection{前缀因果受保护状态}

记 $\Pi_{\mathcal Y}$ 为到 $\mathcal Y$ 的标量裁剪。对于单个物理电芯，PCHP 按下式初始化并更新受保护状态：
\begin{align}
  b_{i1}&=\Pi_{\mathcal Y}(r_{i1}),\\
  b_{it}&=\Pi_{\mathcal Y}\!\left[b_{i,t-1}
  +\alpha\min\{r_{it}-b_{i,t-1},0\}\right],\qquad t>1,
  \label{eq:causal_state_zh}
\end{align}
其中，$\alpha\in(0,1]$ 仅使用源域选择。向上创新被忽略，向下创新按比例 $\alpha$ 同化。因此 $b_{it}\leq b_{i,t-1}$；当 $\alpha<1$ 时，孤立的负向冲击会被衰减。

\begin{lemma}[受保护状态的前缀不变性]
\label{lem:prefix_zh}
对任意 $k\leq T$，在 $(r_{i1},\ldots,r_{ik})$ 上按式~\eqref{eq:causal_state_zh} 计算所得的前 $k$ 个状态，与在任意扩展序列 $(r_{i1},\ldots,r_{iT})$ 上计算所得的前 $k$ 个状态完全相同。
\end{lemma}
\begin{proof}
初始状态只使用 $r_{i1}$，之后每个状态都是 $(b_{i,t-1},r_{it})$ 的确定性函数。对 $t$ 作归纳即可同时得到前缀不变性和关于 $\mathcal F_{it}$ 的可测性。
\end{proof}

\subsection{递归可行的损害预算投影}

令 $p_{i0}=1.3$。在记录 $t$，定义可行区间
\begin{equation}
  I_{it}=\left[
  \max\{0,b_{it}-\delta\},
  \min\{1.3,b_{it}+\delta,p_{i,t-1}\}
  \right],
  \label{eq:feasible_interval_zh}
\end{equation}
并输出
\begin{equation}
  p_{it}=\Pi_{I_{it}}(c_{it}).
  \label{eq:method_zh}
\end{equation}
区间上端点约束当前输出不超过上一条已发布预测，$b_{it}$ 周围的损害管则限制其相对于受保护状态的位移。该投影是在当前可行集合内距离学习候选最近的输出。

\begin{theorem}[递归可行性、前缀因果性与单调性]
\label{thm:pchp_zh}
若 $b_{it}$ 遵循式~\eqref{eq:causal_state_zh}，且 $\delta\geq0$ 为常数，则对所有 $t$，区间 $I_{it}$ 均非空。式~\eqref{eq:method_zh} 的输出具有前缀因果性，属于 $\mathcal Y$，随 $t$ 非增，并满足 $|p_{it}-b_{it}|\leq\delta$。
\end{theorem}
\begin{proof}
当 $t=1$ 时，$p_{i0}=1.3$ 且 $b_{i1}\in\mathcal Y$，故区间非空。假设 $t-1$ 时可行。由 $b_{it}\leq b_{i,t-1}$ 以及 $p_{i,t-1}\geq\max\{0,b_{i,t-1}-\delta\}$，可得 $p_{i,t-1}\geq\max\{0,b_{it}-\delta\}$，所以新的下端点不超过构成上端点的任何一项。投影到 $I_{it}$ 给出取值范围和位移性质，而 $p_{it}\leq p_{i,t-1}$ 给出单调性。所有量只使用已观测前缀，因此输出具有前缀因果性。
\end{proof}

\subsection{损失依赖的精确损害几何}

设预测空间为 $\mathcal P$，结果集合为 $\mathcal Z$，损失函数为 $\ell$，受保护预测为 $b$，预算 $\delta\geq0$。定义结果一致损害域
\begin{equation}
  \mathcal H^{\ell}_{\delta}(b;\mathcal Z)
  =\left\{p\in\mathcal P:
  \sup_{y\in\mathcal Z}\bigl[\ell(p,y)-\ell(b,y)\bigr]\leq\delta
  \right\}.
  \label{eq:generic_harm_region_zh}
\end{equation}
结果集合本身是执行契约的一部分；改变该集合可能改变是否存在可认证的非平凡更新。

\begin{theorem}[度量损失的精确几何]
\label{thm:metric_harm_zh}
设 $(\mathcal M,d)$ 为任意度量空间，并令 $q,y\in\mathcal M$ 上的损失为 $\ell(q,y)=d(q,y)$。对任意 $b,p\in\mathcal M$，有
\begin{equation}
  \sup_{y\in\mathcal M}\left[d(p,y)-d(b,y)\right]=d(p,b),
  \label{eq:metric_identity_zh}
\end{equation}
从而
\begin{equation}
  \mathcal H^{d}_{\delta}(b;\mathcal M)
  =\{p\in\mathcal M:d(p,b)\leq\delta\}.
  \label{eq:metric_ball_zh}
\end{equation}
\end{theorem}
\begin{proof}
由三角不等式，对任意 $y$ 均有 $d(p,y)-d(b,y)\leq d(p,b)$；令 $y=b$ 即可达到等号。求解相应不等式即得式~\eqref{eq:metric_ball_zh}。
\end{proof}

定理~\ref{thm:metric_harm_zh} 是三角不等式在度量空间中的直接推论，本文不把该数学原语作为优先权主张。它的作用在于说明：精确结果一致损害对应一种几何，而不是绝对损失的偶然性质。标量绝对损失是 $\mathcal M=\R$ 且 $d(p,y)=|p-y|$ 的特例。

\begin{theorem}[精确损害域与零预算不可能性]
\label{thm:harm_region_zh}
对固定的 $b\in\R$ 和 $\delta\geq0$，定义
\begin{equation}
  \mathcal H_\delta(b)=\left\{p\in\R:
  \sup_{y\in\R}\bigl(|p-y|-|b-y|\bigr)\leq\delta\right\}.
  \label{eq:harm_region_zh}
\end{equation}
则 $\mathcal H_\delta(b)=[b-\delta,b+\delta]$。因此，
\begin{equation}
  I_{it}=\mathcal H_\delta(b_{it})\cap\mathcal Y
  \cap(-\infty,p_{i,t-1}],
  \label{eq:maximal_interval_zh}
\end{equation}
即 $I_{it}$ 恰好是同时满足物理取值范围、相对于上一条已发布预测的非增约束以及声明最坏绝对损失预算的最大当前记录输出集合。特别地，$\mathcal H_0(b)=\{b\}$：任何非平凡的结果未知更新 $p\neq b$，都会在某个可能真实值上使绝对损失增加，且最坏增量恰为 $|p-b|$。
\end{theorem}
\begin{proof}
由定理~\ref{thm:metric_harm_zh}，式~\eqref{eq:harm_region_zh} 中的定义不等式等价于 $|p-b|\leq\delta$，从而得到区间恒等式。将该精确损害域与 $\mathcal Y=[0,1.3]$ 以及 $p\leq p_{i,t-1}$ 相交，即得式~\eqref{eq:maximal_interval_zh}。令 $\delta=0$ 时只剩 $p=b$；若 $p\neq b$，式~\eqref{eq:metric_identity_zh} 给出一个使损害达到 $|p-b|>0$ 的真实值。
\end{proof}

因此，损害管不是启发式充分条件，而是声明逐点主张的精确可行域。由定理~\ref{thm:harm_region_zh} 与定理~\ref{thm:pchp_zh} 可得逐记录保证
\begin{equation}
  |p_{it}-y_{it}|-|b_{it}-y_{it}|
  \leq |p_{it}-b_{it}|\leq\delta.
  \label{eq:pointwise_harm_zh}
\end{equation}
递归单调约束可能使 $I_{it}$ 比独立损害管更窄。

\begin{theorem}[平方损失的结果范围二分性]
\label{thm:squared_geometry_zh}
设 $L<U$，$b,p\in[L,U]$，并定义 $\ell_2(q,y)=(q-y)^2$。若声明结果集合为 $[L,U]$，则
\begin{equation}
  \sup_{y\in[L,U]}\left[(p-y)^2-(b-y)^2\right]
  =
  \begin{cases}
    (p-b)(p+b-2L), & p\geq b,\\
    (b-p)(2U-p-b), & p<b.
  \end{cases}
  \label{eq:squared_identity_zh}
\end{equation}
因此，精确有限预算域为
\begin{equation}
  \mathcal H^{(2)}_{\delta}(b;[L,U])
  =\left[
  \max\!\left\{L,U-\sqrt{(U-b)^2+\delta}\right\},
  \min\!\left\{U,L+\sqrt{(b-L)^2+\delta}\right\}
  \right].
  \label{eq:squared_region_zh}
\end{equation}
若声明结果集合改为 $\R$，则对所有 $p\neq b$ 均有
\begin{equation}
  \sup_{y\in\R}\left[(p-y)^2-(b-y)^2\right]=+\infty.
  \label{eq:squared_unbounded_zh}
\end{equation}
因此，当结果集合为整个实数轴时，任何非平凡更新都不可能具有有限的结果一致平方损失预算。
\end{theorem}
\begin{proof}
展开损失差可得 $(p-b)(p+b-2y)$。当 $p\geq b$ 时，该仿射函数在 $y=L$ 处取最大值；当 $p<b$ 时，它在 $y=U$ 处取最大值，从而得到式~\eqref{eq:squared_identity_zh}。在 $b$ 的上侧，预算不等式等价于 $(p-L)^2-(b-L)^2\leq\delta$；在下侧，则等价于 $(U-p)^2-(U-b)^2\leq\delta$。分别求解并与 $[L,U]$ 相交，即得式~\eqref{eq:squared_region_zh}。若 $p>b$，展开后的损失差在 $y\to-\infty$ 时发散；若 $p<b$，则在 $y\to+\infty$ 时发散，从而得到式~\eqref{eq:squared_unbounded_zh}。
\end{proof}

\begin{corollary}[有界平方损失下的递归可行性]
\label{cor:squared_recursive_zh}
对固定预算 $\delta\geq0$，定义
\begin{align}
  \underline q^{(2)}_{\delta}(b)
  &=\max\!\left\{L,U-\sqrt{(U-b)^2+\delta}\right\},\\
  \overline q^{(2)}_{\delta}(b)
  &=\min\!\left\{U,L+\sqrt{(b-L)^2+\delta}\right\}.
  \label{eq:squared_bounds_zh}
\end{align}
若 $b_{it}\in[L,U]$ 前缀因果且非增，令 $p_{i0}=U$，并定义
\begin{equation}
  I^{(2)}_{it}
  =\left[
  \underline q^{(2)}_{\delta}(b_{it}),
  \min\!\left\{\overline q^{(2)}_{\delta}(b_{it}),p_{i,t-1}\right\}
  \right].
  \label{eq:squared_interval_zh}
\end{equation}
则对每条记录和任意候选序列，$I^{(2)}_{it}$ 均非空。将 $c_{it}$ 投影到该区间所得输出具有前缀因果性，保持在 $[L,U]$ 内且非增，并满足
\begin{equation}
  (p_{it}-y)^2-(b_{it}-y)^2\leq\delta
  \qquad\text{对所有 }y\in[L,U].
  \label{eq:squared_pointwise_harm_zh}
\end{equation}
\end{corollary}
\begin{proof}
下端点函数 $\underline q^{(2)}_{\delta}(b)$ 关于 $b$ 非减。因此，$b_{it}\leq b_{i,t-1}$ 蕴含 $\underline q^{(2)}_{\delta}(b_{it})\leq\underline q^{(2)}_{\delta}(b_{i,t-1})\leq p_{i,t-1}$。式~\eqref{eq:squared_region_zh} 的独立损害域包含 $b_{it}$，故其下端点也不超过上端点，递归区间因而非空。投影给出取值范围与非增性质，前缀可测性给出因果性，而定理~\ref{thm:squared_geometry_zh} 给出式~\eqref{eq:squared_pointwise_harm_zh}。
\end{proof}

上述二分性解释了本文的实验损失选择，但并不声称平方损失在经验上更差。绝对损失无需限制隐藏真实值即可产生有限精确证书。有界区间平方损失构造只有在 $[L,U]$ 是可信结果契约时才有效；定理~\ref{thm:squared_geometry_zh} 另行证明了结果遍历 $\R$ 时的不可能性。补充材料报告了针对式~\eqref{eq:squared_identity_zh}--\eqref{eq:squared_interval_zh} 的确定性符号、边界、区间外点和随机轨迹证伪测试。

\begin{corollary}[精确非对称决策损失扩展]
\label{cor:asymmetric_zh}
设 $c_{\mathrm{under}},c_{\mathrm{over}}>0$，并定义方向加权绝对损失
\begin{equation}
  \ell_{\mathrm{asym}}(p,y)
  =c_{\mathrm{under}}(y-p)_{+}
  +c_{\mathrm{over}}(p-y)_{+},
  \label{eq:asymmetric_loss_zh}
\end{equation}
其中 $(z)_{+}=\max\{z,0\}$。对任意 $b,p\in\R$，有
\begin{equation}
  \sup_{y\in\R}\!\left[
  \ell_{\mathrm{asym}}(p,y)-\ell_{\mathrm{asym}}(b,y)
  \right]
  =
  \begin{cases}
    c_{\mathrm{over}}(p-b), & p\geq b,\\
    c_{\mathrm{under}}(b-p), & p<b.
  \end{cases}
  \label{eq:asymmetric_identity_zh}
\end{equation}
因此，对决策损失预算 $\eta\geq0$，精确必要充分损害域为
\begin{equation}
  \mathcal H^{\mathrm{asym}}_{\eta}(b)
  =\left[b-\frac{\eta}{c_{\mathrm{under}}},
          b+\frac{\eta}{c_{\mathrm{over}}}\right].
  \label{eq:asymmetric_region_zh}
\end{equation}
将 $I_{it}$ 替换为
\begin{equation}
  I^{\mathrm{asym}}_{it}
  =\left[
  \max\!\left\{0,b_{it}-\frac{\eta}{c_{\mathrm{under}}}\right\},
  \min\!\left\{1.3,b_{it}+\frac{\eta}{c_{\mathrm{over}}},p_{i,t-1}\right\}
  \right],
  \label{eq:asymmetric_interval_zh}
\end{equation}
当代价权重与预算在同一轨迹内保持常数时，仍可保持递归非空、前缀因果性和非增性。令 $c_{\mathrm{under}}=c_{\mathrm{over}}=1$ 且 $\eta=\delta$，即可精确恢复原始 PCHP。
\end{corollary}
\begin{proof}
若 $p\geq b$，损失差的上界为 $c_{\mathrm{over}}(p-b)$，且任取 $y\leq b$ 均可达到该上界；若 $p<b$，上界为 $c_{\mathrm{under}}(b-p)$，且任取 $y\geq b$ 均可达到。分别求解两个方向的不等式即可得到式~\eqref{eq:asymmetric_region_zh}。该损害域的下端点随 $b_{it}$ 非增，因此定理~\ref{thm:pchp_zh} 的归纳论证可不变地用于式~\eqref{eq:asymmetric_interval_zh}。
\end{proof}

方向加权绝对损失是与分位数损失相关的既有预测工具 \citep{koenker1978regression}；推论~\ref{cor:asymmetric_zh} 的贡献并非提出该损失，而是给出其与结果无关的精确损害几何，并将其递归可行地纳入同一在线契约。

\begin{corollary}[时变非对称契约下的精确可行性]
\label{cor:timevarying_zh}
令预算 $\eta_{it}\geq0$ 与代价 $c_{\mathrm{under},it},c_{\mathrm{over},it}>0$ 随 $t$ 变化，并定义
\begin{align}
  L_{it}&=\max\!\left\{0,b_{it}-\frac{\eta_{it}}{c_{\mathrm{under},it}}\right\},
  &
  U_{it}&=\min\!\left\{1.3,b_{it}+\frac{\eta_{it}}{c_{\mathrm{over},it}}\right\},
  \label{eq:timevarying_bounds_zh}\\
  I^{\mathrm{tv}}_{it}&=[L_{it},\min\{U_{it},p_{i,t-1}\}].
  \label{eq:timevarying_interval_zh}
\end{align}
当受保护状态 $b_{it}$ 非增时：
\begin{enumerate}[label=(\roman*),leftmargin=*]
  \item 给定已经实现的前缀，$I^{\mathrm{tv}}_{it}$ 非空当且仅当 $L_{it}\leq p_{i,t-1}$；
  \item 对每一条候选序列均保持递归可行，当且仅当对所有 $t\geq2$ 均有 $L_{it}\leq L_{i,t-1}$；
  \item 给定 $p_{i,t-1}$，精确的当前最小预算为
  \begin{equation}
    \eta^{\min}_{it}
    =c_{\mathrm{under},it}(b_{it}-p_{i,t-1})_{+}.
    \label{eq:minimum_viable_budget_zh}
  \end{equation}
\end{enumerate}
上侧代价日程会改变损害管上侧半径与投影输出，但不决定递归区间是否为空。
\end{corollary}
\begin{proof}
由于 $b_{it}\in[0,1.3]$，独立损害管包含 $b_{it}$，故 $L_{it}\leq U_{it}$。因此，式~\eqref{eq:timevarying_interval_zh} 非空，当且仅当其下端点不超过 $p_{i,t-1}$，从而得到 (i)。若 $L_{it}\leq L_{i,t-1}$，任一可行的上一时刻投影都满足 $p_{i,t-1}\geq L_{i,t-1}\geq L_{it}$，归纳即可证明 (ii) 的充分性。反之，令 $t^{\star}$ 为首个满足 $L_{i,t^{\star}}>L_{i,t^{\star}-1}$ 的索引，并在 $t^{\star}-1$ 之前始终选取不高于下端点的候选；投影将给出 $p_{i,t^{\star}-1}=L_{i,t^{\star}-1}$，于是下一时刻区间为空。因此，更弱的日程条件不能保证对所有候选序列均可行。最后，对 $L_{it}\leq p_{i,t-1}$ 求最小非负 $\eta_{it}$，即可得到式~\eqref{eq:minimum_viable_budget_zh}。下端点不含 $c_{\mathrm{over},it}$，故最后一个结论成立。
\end{proof}

对于预先声明的完整日程，推论~\ref{cor:timevarying_zh} 可以在执行前提供全轨迹证书。如果代价或预算只能按前缀到达，则式~\eqref{eq:minimum_viable_budget_zh} 提供在线可行性诊断：低于该下限的声明预算会使区间为空，系统必须失败关闭，不能静默扩大预算。

\begin{corollary}[精确约束最优性]
\label{cor:optimality_zh}
对每条记录，PCHP 输出是下列问题的唯一解：
\begin{equation}
\begin{aligned}
  p_{it}=\operatorname*{arg\,min}_{p\in\R}\;&\frac{1}{2}(p-c_{it})^2\\
  \text{约束为}\quad
  &p\in\mathcal Y,\qquad p\leq p_{i,t-1},\\
  &\sup_{y\in\R}\bigl(|p-y|-|b_{it}-y|\bigr)\leq\delta.
\end{aligned}
\label{eq:constrained_optimality_zh}
\end{equation}
\end{corollary}
\begin{proof}
由定理~\ref{thm:harm_region_zh}，式~\eqref{eq:constrained_optimality_zh} 的可行集恰为 $I_{it}$；定理~\ref{thm:pchp_zh} 保证该区间非空。严格凸目标因而具有唯一极小点 $\Pi_{I_{it}}(c_{it})$，即式~\eqref{eq:method_zh}。
\end{proof}

\begin{corollary}[聚合继承]
对任意非负权重 $w_k$ 和预算 $\delta_k$，均有
\begin{equation}
  \sum_k w_k\left(|p_k-y_k|-|b_k-y_k|\right)
  \leq\sum_k w_k\delta_k.
  \label{eq:aggregate_zh}
\end{equation}
若 $\sum_k w_k=1$，则加权 MAE 的增加不超过预算的加权平均。
\end{corollary}

式~\eqref{eq:aggregate_zh} 无需独立性假设。经验效用的不确定性仍是独立问题，必须遵守数据域—电芯—循环的层级结构。实验采用的常数预算是递归可行的特例；推论~\ref{cor:timevarying_zh} 给出了时变契约所需的精确附加条件。

\subsection{计算结构}

两个冻结回归器生成当前记录预测后，PCHP 只需保存 $(b_{i,t-1},p_{i,t-1})$，并执行常数次标量运算。因此，模型推理之后，每条记录只需 $O(1)$ 状态与 $O(1)$ 时间。有界标量平方损失扩展仍保持相同复杂度，因为式~\eqref{eq:squared_region_zh} 的两个端点均为闭式。审计链记录 $r_{it}$、$c_{it}$、$b_{it}$、$I_{it}$、$p_{it}$、$\alpha$ 与 $\delta$。与回顾性等距回归不同，在线阶段无须对完整轨迹排序，也无须在未来记录到达后重新计算历史输出。

\section{实验设计}

\subsection{开发数据域与嵌套模型选择}

开发数据由 BatteryLife 处理数据快照（v11）、BatteryML、Battery Archive，以及 CALCE、HNEI、HUST、MATR、MICH、MICH\_EXP、RWTH、SNL、UL\_PUR、SDU 与 XJTU 各数据域对应的原始实验来源标准化整理 \citep{zhang2024batteryml,tan2025batterylife,batterylife2026processed,batteryarchive,xing2013ensemble,devie2018variability,ma2022personalized,severson2019cyclelife,attia2020closedloop,weng2021formation,li2021oneshot,preger2020degradation,juarezrobles2020safety,wang2025sorting,wang2024pinn}。去除前五次参照窗口后，分析语料包括 $12$ 个完整数据域、$586$ 个物理电芯和 $601{,}932$ 条记录，见表~\ref{tab:data_zh}。这 $12$ 个数据域均已在方法开发阶段打开，因此不被表述为独立外部确认。

\begin{table}[t]
\centering
\caption{前五次参照窗口后的开发语料。记录嵌套于物理电芯，电芯嵌套于完整数据域。}
\label{tab:data_zh}
\begin{tabular}{lrr@{\qquad}lrr}
\toprule
数据域 & 电芯数 & 记录数 & 数据域 & 电芯数 & 记录数 \\
\midrule
CALB & 27 & 2,629 & MICH\_EXP & 18 & 6,455 \\
CALCE & 13 & 14,233 & RWTH & 48 & 105,140 \\
HNEI & 14 & 15,085 & SDU & 86 & 48,716 \\
HUST & 77 & 145,735 & SNL & 61 & 96,092 \\
MATR & 169 & 139,148 & UL\_PUR & 10 & 2,195 \\
MICH & 40 & 19,681 & XJTU & 23 & 6,823 \\
\midrule
\multicolumn{4}{r}{合计} & 586 & 601,932 \\
\bottomrule
\end{tabular}
\end{table}

外层评估每次留出一个完整目标域，两个回归器只在其余十一个数据域上拟合。在这些源域内部，对每一个候选值
\[
\alpha\in\{1,0.5,0.2,0.1,0.05,0.03,0.02,0.015,0.01,0.005\}
\]
再次逐域留出，并在其余十个数据域上拟合模型。以内层域等权电芯宏平均 MAE 最小的 $\alpha$ 作为外层折选择；当差异不超过 $10^{-12}$ 时选择较大的 $\alpha$。外层目标域标签从不参与这一选择。全部因果分析均固定使用 $\delta=0.01$ 个 SOH 单位的损害预算。由于该预算形成于开发阶段，其工程决策含义必须由实际部署者给出，不能从外部压力测试中反推。

对每个电芯，首先在其全部参照窗口后记录上计算 MAE；随后在域内对电芯等权平均，最后对 $12$ 个数据域等权平均。主要开发估计量是 PCHP 与其 $\alpha$ 匹配受保护因果基线之间的域等权配对差。该层级避免大型数据档案或长寿命电芯支配跨域结论。

\subsection{损失几何证伪协议}

在将损失扩展写入论文之前，本文冻结了符号恒等式、有限区间闭式、区间外点失败检查、递归可行门、$10^{-12}$ 容差和随机种子 $20{,}260{,}802$。该确定性测试不读取任何电池结果。测试在 $200{,}000$ 个随机配置中比较有界平方损失闭式与端点极大化结果，检查可行域两侧边界及其相邻区间外点，并在 $2{,}000$ 条长度为 $128$ 的随机非增受保护轨迹上投影任意候选。另行构造两个反例，分别沿正确的无穷方向移动隐藏结果，以检验非平凡向上和向下更新的平方损失发散。任何公式不一致、可行点预算违反、区间外点仍可行、递归区间为空、输出上升或未按预期发散，均要求在不调整容差的条件下拒绝该扩展。

\subsection{机制证伪与区分性控制}

嵌套评分前，最小机制证伪只使用不含标签的开发预测。在轨迹比例 $0.25$、$0.50$ 和 $0.75$ 处，本文比较仅用当前前缀计算的输出与完整未来轨迹到达后重新计算的相同历史位置。若变化超过 $10^{-4}$ 个 SOH 单位，则记为一次回顾性修订。因果机制必须保持最大前缀差为 $0$；完整轨迹等距回归则应在至少九个数据域中修改不少于一半电芯的历史前缀。第二项干预在单条受保护原始预测上加入 $+0.05$ 或 $-0.05$ 的创新。预期特征为：正向创新被完全拒绝，负向创新即时按 $\alpha$ 同化，并且相对于 $\alpha=1$，$\alpha=0.02$ 应在至少九个数据域中对不少于 $95\%$ 的干预降低累计扰动。

主要比较对象是使用同一源域选择 $\alpha$ 生成的因果受保护基线。回顾性等距损害预算方法仅作为信息泄漏审计保留：它能够访问完整目标轨迹，因此不满足本文的部署信息边界。最强替代解释是，PCHP 可能只是把受保护状态在同一 $\delta$ 损害管内做常数平移。为检验该解释，候选无关控制对每个外层折仅使用十一个源域，联合选择 $\alpha$ 和 $\{-0.01,-0.009,\ldots,0.01\}$ 中的常数偏移。该控制与 PCHP 具有相同的取值范围、单调性和损害预算约束。

在开发结果已经打开之后，本文进一步协议锁定了一项回顾性预算路径敏感性分析，用于检验上述区分是否依赖于单一设定 $\delta=0.01$。计算前冻结
\[
\delta\in\{0,0.0025,0.005,0.0075,0.01,0.015,0.02,0.03\},
\]
原始 $\alpha$ 网格、外层数据域划分、唯一的积分主要估计量、bootstrap 随机种子 $20{,}260{,}811$ 以及判定门槛。对每个外层数据域、预算和 $\alpha$，候选无关比较器先计算源域残差的加权中位数，即使内层域等权电芯宏平均绝对误差最小的精确常数偏移，再将其截断到 $[-\delta,\delta]$；随后仍只用同一组十一个源域选择 $\alpha$。两项选择均不使用外层目标域标签。零预算组必须严格复现受保护状态，所有预算组均须通过声明的位移、遗憾、取值范围和单调性检查。该分析虽在执行前锁定协议，但仍属于回顾性分析，而非新的确认实验。

本文还构造了更强的\emph{无保护}因果比较器，用于度量损害管的机会成本。对每个外层目标域，仍采用仅用源域的双重留出设计，为下式选择 $\alpha$：
\[
q_{i1}=\Pi_{[0,1.3]}(c_{i1}),\qquad
q_{it}=\Pi_{[0,1.3]}\!\left(q_{i,t-1}+\alpha\min(c_{it}-q_{i,t-1},0)\right).
\]
该比较器具有前缀因果性、取值有界且轨迹非增，但不受保护基线周围损害管的约束。十二个外层折均仅根据源域选择 $\alpha=0.05$。因此，它是运行上较强的精度边界，而不是安全性匹配的基线。

\subsection{结果隔离的 BaSyTec 外部确认}

拜罗伊特退化路径数据集包含 $48$ 个商用锂离子电芯，在 $24$ 组温度—倍率组合下进行 $70$ 个老化循环，每个工况名义上有两个重复电芯 \citep{ramasubramanian2025degradation,roder2025degradationdata}。评分前，本文固化了 Zenodo 公开记录、全部 $48$ 个电芯压缩包及其官方校验和。第一版解析协议在读取任何容量数值之前闭门失败：原始导出只包含累计电量字段 \texttt{Ah[Ah]}，没有逐循环容量字段；该失败被完整保留。在查看任何累计电量数值之前，F0001 被永久指定为解析开发电芯并排除于确认。随后对其余 $47$ 个压缩包执行只读表头、零数据行的结构普查，在预测前冻结两类分隔符和四种表头签名。

对循环 $c$ 内相邻行 $j-1,j$，记累计电量为 $A_j$、电流为 $I_j$，冻结的循环放电容量恢复式为
\begin{equation}
Q_c=\sum_{j=2}^{n_c}\mathbf{1}\!\left(I_j<-10^{-5}\,\mathrm{A}\right)
\max(A_{j-1}-A_j,0).
\label{eq:basytec_capacity_zh}
\end{equation}
标签盲阶段只选择循环标识以及充电侧时间、电流和电压；数值小时乘以 $3600$ 转换为秒，随后提取固定 $600\,\mathrm{s}$ 充电窗口。在计算式~\eqref{eq:basytec_capacity_zh} 之前，程序为 $47$ 个电芯生成并哈希冻结 $3{,}101$ 条参照后预测。评分程序随后按式~\eqref{eq:soh_zh} 构造目标。电芯必须满足前五循环中位容量位于 $[0.08,0.16]\,\mathrm{Ah}$，并至少具有十条参照后标签。PCHP 固定使用 $\alpha=0.01$、$\delta=0.01$；无保护边界使用源域选择的 $\alpha=0.05$。

判定门槛要求至少 $36$ 个合格电芯、电芯等权配对 MAE 差为负、至少三分之二非并列电芯改善、电芯 bootstrap $95\%$ 区间上界小于 $0$，且最大逐点位移、最大观察绝对损失遗憾和最大电芯宏平均损害均不超过 $0.01$（容差 $10^{-12}$）。当前工作区记录支持“预测哈希先于容量访问”，但无法独立排除团队在工作区外的历史暴露。因此，本文称其为结果隔离、协议冻结的公开数据确认，而不是前瞻性实验室试验。

\subsection{协议冻结的 NASA 一次性压力测试}

NASA PCoE Battery Data Set~5 依据官方元数据选定，并在下载压缩包之前完成冻结 \citep{saha2007nasa}。协议固定了官方压缩包身份、递归发现的全部 MATLAB 文件纳入规则、充放电配对、$600\,\mathrm{s}$ 充电窗口、前五次参照、允许读取的特征字段、$\alpha=0.01$、$\delta=0.01$、最小电芯记录数、统计估计量、判定门槛，以及预测与评分程序哈希。去除重复身份后，压缩包中的 $38$ 个 MATLAB 文件对应 $34$ 个物理电池。

第一阶段只访问循环类型、时间，以及充电记录中的电压、电流、温度和时间数组；在不读取放电容量的条件下，为 $34$ 个电池生成并哈希冻结 $2{,}598$ 条参照后预测。第二阶段由独立程序读取冻结的放电定位信息和标量 \texttt{Capacity} 字段，按式~\eqref{eq:soh_zh} 构造结果，并对冻结预测评分一次。每个电池至少需要十条有效参照后记录。冻结结果契约最终得到 $33$ 个电池和 $2{,}556$ 条评分记录；按预先声明的规则，$44$ 条非标量、非有限或非正容量记录被拒绝。结果打开后，不允许改变电池清单、参照定义、模型、参数或终点。

外部门槛要求电池等权配对 MAE 差为负、至少 $75\%$ 的合格电池改善、最大电池宏平均损害不超过 $0.01$，且全部确定性证书通过。该数据只代表一个数据仓库身份和一个实验室项目；物理电池可以作为数据集内重复单位，但不能被解释为相互独立的实验室。在当前可审计项目谱系中，冻结预测产物先于对 Dataset~5 放电容量的访问；但无法独立排除研究团队在该工作区之外的历史结果暴露。因此，本文将其表述为协议冻结的公开数据压力测试，而不是前瞻性外部确认或前瞻性实验室试验。

\subsection{统计分析}

开发效应以完整数据域为主要分析单位，BaSyTec 与 NASA 数据集内效应以物理电芯或电池为分析单位；重复记录均视为嵌套测量，不计作独立重复。开发主效应与候选控制分别使用固定随机种子 $20{,}260{,}802$ 和 $20{,}260{,}801$，对 $12$ 个域配对差执行 $100{,}000$ 次百分位 bootstrap。预算路径预先指定的逐域估计量为
\begin{equation}
A_d=\frac{1}{0.03}\int_{0}^{0.03}
\left[\operatorname{MAE}^{\mathrm{PCHP}}_d(\delta)-
\operatorname{MAE}^{\mathrm{shift}}_d(\delta)\right],\mathrm d\delta,
\label{eq:budget_auc_zh}
\end{equation}
并在冻结网格上用梯形法计算。主要报告量是 $12$ 个 $A_d$ 的域等权均值及其 $100{,}000$ 次完整域 bootstrap 区间；各预算点的区间只是同一条路径的描述性切片，不构成多项独立确认主张。鉴于 MICH 与 MICH\_EXP 来自同一实验项目，另行实施保守的事后敏感性分析：先合并二者，再按实验来源组等权聚合与重采样，并依次删除每个来源组重算均值。BaSyTec 使用随机种子 $20{,}260{,}808$，对物理电芯执行 $100{,}000$ 次配对 bootstrap；另行实施事后工况敏感性，将同一温度—倍率下的重复电芯先平均，再对完整工况重采样。NASA 使用随机种子 $20{,}260{,}801$，对物理电池执行 $100{,}000$ 次配对 bootstrap。双侧精确符号检验排除平局，只作为描述性统计量而非确认门槛。由于全部开发域均已打开，其区间用于概括已观察数据域之间的异质性，不构成对预先指定总体的确认性覆盖。两个公开数据流程采用冻结的多准则判定门槛，而非经显著性水平校准的假设检验；本文不提出族错误率或多重校正推断主张。每个区间只对应所述比较；本文不执行循环层级检验，也不基于合并记录构造置信区间。确定性证书属于逐记录数学性质，不依赖这些重采样假设。

\section{结果}

\subsection{损失依赖几何通过符号与递归证伪}

冻结的 $27$ 个聚合检查全部通过。在 $200{,}000$ 个有界平方损失配置中，端点极大化结果与式~\eqref{eq:squared_identity_zh} 的最大绝对差为 $1.33\times10^{-15}$；有效下边界和上边界相对于声明预算的最大差分别为 $8.88\times10^{-16}$ 与 $1.55\times10^{-15}$。在 $256{,}000$ 条递归记录中，没有出现空区间，最大输出增量为 $0$，最大预算超额为 $7.77\times10^{-16}$。两个无界结果构造均沿式~\eqref{eq:squared_unbounded_zh} 预测的方向发散。这些测试不替代理论证明，其作用是排除实现中的符号、端点、截断与递归错误。

\subsection{因果算子避免未来修订并呈现预期冲击响应}

最小证伪的两个门槛均通过。完整轨迹等距回归在全部 $12$ 个数据域中修改了超过一半的电芯前缀，而因果前缀与其位于任意未来扩展中的同一前缀之间，最大差值严格为 $0$。施加 $+0.05$ 的受控创新后，单向受保护状态的最大扰动为 $0$；施加 $-0.05$ 创新后，在已保存干预明细表中，即时响应相对于理论值 $\alpha\times0.05$ 的最大误差不超过 $2.00\times10^{-16}$，处于浮点容差范围内。相对于 $\alpha=1$，$\alpha=0.02$ 在每个数据域中均对至少 $95\%$ 的干预降低了累计冲击扰动。图~\ref{fig:causality_zh} 表明，该机制真正改变了信息边界，而不是把离线平滑重新命名。

\begin{figure}[t]
  \centering
  \includegraphics[width=\textwidth]{fig3_prefix_causality_zh.pdf}
  \caption{覆盖 $12$ 个开发数据域和 $586$ 个物理电芯的机制证伪。\textbf{A}，未来记录追加后，历史离线输出变化超过 $10^{-4}$ 的电芯比例；PCHP 的最大前缀差为 $0$，虚线表示预先声明的半数电芯门槛。\textbf{B}，$\alpha=1$ 与 $\alpha=0.02$ 下累计冲击扰动的逐域中位数配对，横轴采用对数尺度；每个数据域中的全部干预均得到改善。记录和干预嵌套于电芯，完整数据域是跨域评估单位。}
  \label{fig:causality_zh}
\end{figure}

\subsection{嵌套源域选择在损害预算内提高跨域效用}

在严格嵌套留一数据域评估中，PCHP 将域等权电芯宏平均 MAE 从因果受保护基线的 $0.07524$ 降至 $0.07038$，相对降低 $6.46\%$。配对差为 $-0.00486$；按完整数据域聚类重采样得到的 $95\%$ 百分位自助区间为 $[-0.00671,-0.00301]$。其中 $11$ 个数据域改善，无平局，仅一个数据域恶化。最大域层级损害为 $0.00109$，明显小于 $0.01$ 的逐记录预算。图~\ref{fig:domain_effects_zh} 展示全部数据域，避免把 $601{,}932$ 条重复记录误作独立跨域证据。

\begin{figure}[t]
  \centering
  \includegraphics[width=0.92\textwidth]{fig2_domain_effects_zh.pdf}
  \caption{全部 $12$ 个外层目标域的 PCHP 减基线配对电芯宏平均 MAE。每个点代表一个经电芯等权聚合的完整数据域，负值表示 PCHP 更优。菱形及水平区间分别为域等权均值和对完整数据域执行 $100{,}000$ 次百分位 bootstrap 得到的 $95\%$ 区间。}
  \label{fig:domain_effects_zh}
\end{figure}

因果方法的点估计也低于回顾性损害预算方法，即 $0.07038$ 对 $0.07232$。两者配对差为 $-0.00195$，$95\%$ 区间为 $[-0.01073,0.00606]$。该比较不能证明因果版本在统计上更优；它只说明，在已打开的开发语料中，落实可部署信息边界没有呈现明确的域层级精度代价。

\subsection{学习候选提供了超越常数安全偏移的信息}

候选无关控制的域等权电芯宏平均 MAE 为 $0.07515$，PCHP 为 $0.07038$，相对降低 $6.35\%$。以 PCHP 减控制表示，配对差为 $-0.00477$，域 bootstrap $95\%$ 区间为 $[-0.00808,-0.00159]$。PCHP 在 $9$ 个数据域改善，在 $3$ 个数据域恶化，相对于该控制的最大域损害为 $0.00424$。

协议锁定的回顾性预算路径把控制从离散网格近似升级为精确源域最优常数偏移，并把比较扩展到零至三个 SOH 百分点。七个正预算的域等权均值均支持 PCHP；差值从 $\delta=0.0025$ 时的 $-0.00189$ 扩大至 $\delta=0.03$ 时的 $-0.01205$，且每个预算点的描述性区间均低于零。式~\eqref{eq:budget_auc_zh} 的主要预算归一化 AUC 差值为 $-0.00694$，域 bootstrap 区间为 $[-0.01108,-0.00296]$，逐域胜/平/负为 $11/0/1$。零预算恒等性以及位移、遗憾、取值范围和单调性门槛全部通过。独立哈希固定审计保留了原始验证器的失败记录，同时确认其唯一失败锚点是在 $10^{-12}$ 容差下，把高精度数值 $0.0703775121469438$ 与仅保留十位小数的 JSON 展示值比较所致；权威 CSV 汇总、全部 $12$ 个逐域数值和仅用源域的选择清单均与历史运行一致。因此，图~\ref{fig:candidate_control_zh} 削弱了“结论仅由单一预算下固定偏差修正造成”的解释，但不改变证据边界：开发结果此前已经打开，该路径属于稳健性证据，而不是外部确认。

来源组敏感性分析不依赖于把 MICH 与 MICH\_EXP 视为两个独立实验项目。合并二者后，开发主效应的来源组等权配对差为 $-0.00469$，$95\%$ bootstrap 区间为 $[-0.00667,-0.00274]$，胜/平/负为 $10/0/1$；候选控制配对差为 $-0.00467$，区间为 $[-0.00825,-0.00125]$，胜/平/负为 $8/0/3$。对主效应和候选控制而言，依次删除任意一个来源组后的均值分别始终保持在 $[-0.00526,-0.00417]$ 与 $[-0.00556,-0.00347]$。该分析是处理共同来源的事后敏感性检验，不替代冻结的完整数据域估计量。

\begin{figure}[t]
  \centering
  \includegraphics[width=\textwidth]{fig4_candidate_control_zh.pdf}
  \caption{候选信息的协议锁定回顾性预算路径控制。\textbf{A}，冻结预算网格上 PCHP 减精确安全偏移的域等权电芯宏平均 MAE；负值表示 PCHP 更优，纵向区间为描述性域 bootstrap 区间，菱形标出正文设定 $\delta=0.01$。对每个预算和外层折，控制不使用候选特征，只利用其余十一个源域选择 $\alpha$ 以及精确源域最优常数安全偏移。\textbf{B}，式~\eqref{eq:budget_auc_zh} 所定义的 $12$ 个逐域预算归一化 AUC 差值；菱形和区间为域等权均值及对完整数据域执行 $100{,}000$ 次重采样所得的百分位 bootstrap $95\%$ 区间。该协议在开发结果已经打开后才锁定，因此图中证据属于回顾性机制敏感性分析，而非独立确认。}
  \label{fig:candidate_control_zh}
\end{figure}

\subsection{强无保护比较器揭示安全—精度机会成本}

源域调优的无保护因果比较器将域等权电芯宏平均 MAE 从 $0.07524$ 降至 $0.05470$。它相对于受保护基线的配对差为 $-0.02054$，域 bootstrap $95\%$ 区间为 $[-0.03364,-0.00675]$，胜/平/负为 $10/0/2$。PCHP 保留了该增益的 $23.7\%$：其 MAE 比无保护比较器高 $0.01568$，区间为 $[0.00350,0.02725]$，胜/平/负为 $3/0/9$。合并 MICH 与 MICH\_EXP 后，效用保留率为 $24.9\%$，定性结论不变。

这种额外精度并非安全性匹配。无保护比较器在全部 $12$ 个域、全部 $586$ 个电芯和全部 $524{,}165$ 条评分记录上超过 $\delta=0.01$ 位移预算，最大位移为 $0.35680$；PCHP 没有预算违规。因此，该比较否定了精度支配主张：PCHP 并不在精度上击败最强因果候选。其贡献是在声明的最坏损害区域内，保留可测量的一部分可用效用。

\subsection{冻结 BaSyTec 确认通过，但解释边界十分明确}

BaSyTec 门槛返回 \texttt{RETAIN}。在 $47$ 个确认电芯中，$45$ 个满足冻结参照容量范围，形成 $2{,}969$ 条评分循环。电芯等权 MAE 从受保护因果基线的 $0.17496$ 降至 PCHP 的 $0.16510$。配对差为 $-0.00986$，配对电芯 bootstrap $95\%$ 区间为 $[-0.01000,-0.00959]$；全部 $45$ 个合格电芯均改善。独立重算在 $10^{-12}$ 内复现全部电芯聚合量；事后工况等权敏感性在所覆盖的 $23$ 个温度—倍率工况中均保持负值，区间为 $[-0.01000,-0.00958]$。PCHP 的最大位移和最大观察绝对损失遗憾在浮点容差内均为 $0.01$，最大电芯宏平均差为 $-0.00465$。

$0\,^{\circ}\mathrm{C}$、$1.5C$ 工况下的两个电芯均被预先设定的 $0.08\,\mathrm{Ah}$ 参照下限排除：F0009 与 F0048 的前五循环中位容量分别为 $0.06862$ 和 $0.06935\,\mathrm{Ah}$。明确标注为事后的全样本敏感性在 $47$ 个电芯和 $24$ 个工况上仍保持结论：配对差为 $-0.00986$，区间为 $[-0.01000,-0.00961]$，$47/47$ 个电芯改善。该敏感性表明排除规则没有制造效应，但不能替代冻结的 $45$ 电芯估计量。

同一确认也说明经验主张为何必须收窄。源域调优无保护比较器的 MAE 为 $0.06172$，但最大位移达到 $0.28508$，并在全部合格电芯和全部 $2{,}969$ 条记录上违反预算；PCHP 仅保留其精度增益的 $8.7\%$。此外，受保护基线对每条评分记录都低估，PCHP 在 $99.1\%$ 的记录上达到向上 $0.01$ 的管边界；事后固定上移 $0.01$ 的 MAE 为 $0.16496$，比 PCHP 低 $0.00014$。因此，图~\ref{fig:basytec_zh} 验证了冻结损害—效用契约并量化其机会成本；该外部数据域本身不能区分自适应 PCHP 与简单安全偏差修正。

\begin{figure}[t]
  \centering
  \includegraphics[width=\textwidth]{fig6_basytec_confirmation_zh.pdf}
  \caption{结果隔离的 BaSyTec 确认与安全—精度边界。\textbf{A}，全部 $45$ 个确认合格电芯的 PCHP 减基线配对 MAE，负值表示 PCHP 更优；菱形和区间为电芯等权均值及对物理电芯执行 $100{,}000$ 次配对重采样所得的 $95\%$ bootstrap 区间。\textbf{B}，电芯等权 MAE 与相对受保护基线的最大逐点位移。PCHP 位于声明的 $\delta=0.01$ 区域内；固定上移 $0.01$ 是事后诊断，而源域调优无保护比较器具有前缀因果性但违反全部合格电芯的预算。循环嵌套于电芯。}
  \label{fig:basytec_zh}
\end{figure}

\subsection{冻结 NASA 压力测试保留相对效用并暴露目标不兼容}

NASA 一次性门槛返回 \texttt{RETAIN}。在 $33$ 个合格物理电池中，电池等权 MAE 从因果受保护基线的 $0.73379$ 降至 PCHP 的 $0.72791$。

配对差为 $-0.00588$，电池 bootstrap $95\%$ 区间为 $[-0.00784,-0.00373]$；$27$ 个电池改善，$6$ 个恶化。最大电池层级损害为 $0.00849$。双侧精确符号检验值为 $0.000324$，由于它不是冻结门槛，仅作描述性报告。全部预测均满足 $|p_{it}-b_{it}|\leq0.01$；在浮点容差内，最大观察绝对损失增量为 $0.01$，所有受保护轨迹和最终轨迹均非增。

绝对误差不能支持“NASA 高精度 SOH”主张。一次性结果解封后的冻结标签审计显示，一些放电容量无法在前五次参照下保持协议可比。B0041 的参照容量为 $0.05584\,\mathrm{Ah}$，最大归一化值达到 $21.77$；B0039 的参照容量为 $0.47114\,\mathrm{Ah}$，最大归一化值达到 $3.76$。全部已释放标签的第 $95$ 百分位为 $1.347$，最大值为 $21.77$，而模型输出已经冻结在 $[0,1.3]$。因此，图~\ref{fig:nasa_zh} 将相对效应与协议边界同时呈现。若在结果打开后删除这些电池或重定义参照，将破坏冻结实验，本文没有这样处理。

\begin{figure}[t]
  \centering
  \includegraphics[width=\textwidth]{fig5_nasa_stress_zh.pdf}
  \caption{协议冻结的 NASA 压力测试。\textbf{A}，全部 $33$ 个合格物理电池的 PCHP 减基线配对 MAE，负值表示 PCHP 更优；菱形及区间为电池等权均值和执行 $100{,}000$ 次配对电池重采样所得的 bootstrap $95\%$ 区间。\textbf{B}，按电池展示的已释放归一化容量标签，纵轴采用对数尺度，并标出冻结模型范围 $[0,1.3]$。该图同时说明相对安全结果为何有效，以及高精度绝对 SOH 主张为何不成立。$2{,}556$ 条记录均为嵌套于电池的重复测量。}
  \label{fig:nasa_zh}
\end{figure}

\begin{table}[t]
\centering
\caption{关键配对效应。差值定义为比较方法 MAE 减参照方法 MAE；区间对所述独立单位重采样，而非对单条记录重采样。无保护一行是精度边界，并违反损害预算。}
\label{tab:primary_zh}
\small
\begin{tabularx}{\textwidth}{>{\raggedright\arraybackslash}Xrrrrl}
\toprule
证据层级与比较 & 独立单位数 & 参照 MAE & 比较方法 MAE & 配对差 & $95\%$ 区间 \\
\midrule
嵌套开发主效应 & 12 个域 & 0.07524 & 0.07038 & $-0.00486$ & $[-0.00671,-0.00301]$ \\
PCHP 对候选无关安全偏移 & 12 个域 & 0.07515 & 0.07038 & $-0.00477$ & $[-0.00808,-0.00159]$ \\
无保护候选对受保护基线 & 12 个域 & 0.07524 & 0.05470 & $-0.02054$ & $[-0.03364,-0.00675]$ \\
BaSyTec PCHP 对基线 & 45 个电芯 & 0.17496 & 0.16510 & $-0.00986$ & $[-0.01000,-0.00959]$ \\
NASA 一次性压力测试 & 33 个电池 & 0.73379 & 0.72791 & $-0.00588$ & $[-0.00784,-0.00373]$ \\
\bottomrule
\end{tabularx}
\end{table}

\section{讨论}

\subsection{方法改变了什么能力}

PCHP 改变的是跨域电池预测的运行契约。学习模型仍负责提出有用更新，但不再单方面决定已经部署的健康记录。受保护因果状态明确了需要保护的预测，递归可行区间则精确定义无目标标签候选可以移动多远。推论~\ref{cor:optimality_zh} 表明，最终输出不是该区间内的任意一点，而是在完整声明约束下对当前学习候选保持最大保真度的唯一更新。最终保证既不依赖未知结果，又可以在线执行。它不同于回顾性平滑：后者即使不读取标签，仍可能违反未来信息边界；它也不同于无约束目标适应：后者在未知数据域中的增益可能转为负值。

定理~\ref{thm:harm_region_zh} 还排除了对“安全更新”的误读：在不对隐藏真实值作任何假设时，任何方法只要移动点预测，就不可能相对于基线对所有可能真实值都保持不劣。因此，正预算并不是为了证明方便而引入的松弛，而是允许任何非平凡结果未知更新所必须支付的代价。PCHP 将这一代价显式化，并且只在精确损害域与在线轨迹约束的交集中使用它。

候选控制结果使该方法不能被简化为保守裁剪。如果源域选择的常数偏移已经能够匹配 PCHP，则学习候选并无必要，论文贡献也将退化为通用安全管。实际结果表明，在安全偏移和同化率均得到匹配后，投运参照、记录特定的信息仍提供额外效用。积分预算路径还说明，这一区分并不是单一正文预算造成的偶然现象：在每个冻结预算下，比较器都是常数安全偏移类中的源域最优解，而域等权优势仍贯穿全部正预算网格。该结果不识别电化学因果机理；它识别的是候选特定信息在受控输出契约内的运行价值。预算路径也不能标定工程容差，因为它仍建立在已经打开的开发域上。并且，这种区分并非在每个域都成立：BaSyTec 存在近乎一致的基线低估，事后固定正向平移的精度略高于 PCHP。因此，自适应价值来自嵌套开发控制及其协议锁定敏感性分析，而不是每个外部数据域。

\subsection{BaSyTec 确认与安全—精度边界的含义}

BaSyTec 实验检验的是：在新的实验室、电芯类型和 $24$ 工况应力网格上，容量访问之前冻结的预测产物能否在保持全部声明损害检查的同时改善受保护基线。该检验通过；事后全工况敏感性进一步表明，预设的低参照排除没有制造效应方向。由于评分电芯的归一化容量保持在合理 SOH 区域，该结果提供了比 NASA 压力测试更强的效用证据。

同一实验也阻止了夸大解释。几乎全部 PCHP 输出都达到损害管的正向边界，而固定正向平移的精度略高。因此，该结果支持安全—相对效用契约的迁移，却不能在外部识别自适应候选机制。强无保护比较器则揭示另一条边界：更高精度确实可得，但代价是每个电芯都离开声明损害管。这不是被保证掩盖的缺陷，而是保证使其可量化的工程权衡。仅保留 $8.7\%$ 效用是否可接受，取决于当前公开数据无法识别的实际决策代价。

\subsection{NASA 结果的正确含义}

NASA 实验检验的是：在显著数据源偏移后，完全冻结的因果更新能否保留相对效用和确定性损害控制。该检验通过。它不检验冻结源模型能否在所有 NASA 放电制度下实现可部署的绝对 SOH 精度。异常增大的归一化容量表明，当前五次放电与后续记录采用明显不同的运行条件时，既定目标定义会失去可比性。因此，合理推断必须有边界：在严重协议偏移下，PCHP 仍在声明预算内改善了受保护基线，但该基线和标签契约仍不足以支持绝对部署。

这也说明确定性主张与经验主张必须分开。即使基线误差很大，式~\eqref{eq:pointwise_harm_zh} 仍然成立，因为它在同一记录上比较两个预测。绝对准确性、校准性与运营效用则需要协议兼容的目标和具有代表性的部署数据；任何输出空间定理都不能替代这些条件。

\subsection{损害预算的选择}

预算 $\delta$ 以 SOH 单位表示，应当对应修改受保护记录所带来的后果。较小预算更紧密地保护基线，但限制候选效用；较大预算允许更多收益，也允许更大的最坏绝对损失增量。本文的 $\delta=0.01$ 是冻结研究设定，而不是普适工程容差。当低估与高估的后果不同时，推论~\ref{cor:asymmetric_zh} 使用相应正代价权重和以决策损失为单位的预算 $\eta$，从而形成不同的向下与向上允许半径。由于独立区间只由代价比及与之共同缩放的预算决定，运营者必须通过维护或质保后果明确并记录该契约，不能从本文的精度实验反推代价。在时变契约下，式~\eqref{eq:minimum_viable_budget_zh} 还能报告保持已发布前缀所需的最小预算；该式用于诊断不可行声明，并不标定实际运营代价。阈值型或非线性损失仍需重新推导可行域并给出证明。

\subsection{统计解释}

确定性证书逐记录成立，不要求独立抽样；经验效用的不确定性具有不同适用范围。开发阶段的 $12$ 个完整数据域支持跨域效应，$601{,}932$ 条记录并不构成同等数量的独立迁移；来源组敏感性进一步表明，结果并非由同一实验项目的两个子集被分别计数所驱动。BaSyTec 的 $45$ 个合格电芯支持冻结的数据集内效应，所覆盖的 $23$ 个工况支持事后工况聚类敏感性；全 $47$ 电芯、$24$ 工况分析明确属于事后分析。NASA 的 $33$ 个电池也只支持项目内部重复。两个公开流程都不产生独立实验室重复。因此，本文依照相应分析对完整数据域、保守来源组、电芯、工况或电池重采样。所报告区间刻画已观察单位集合上的异质性，不能证明对所有化学体系、实验室、温度和工况均成立。

\section{局限性与失效模式}

第一，保证相对于指定受保护因果状态成立。PCHP 不能诊断失准基线、修复不兼容标签、建立校准性或保证电化学安全。第二，定理~\ref{thm:metric_harm_zh} 覆盖度量损失，而定理~\ref{thm:squared_geometry_zh} 给出有界区间上的平方损失构造；若结果界不真实或不可核验，则不能使用该保证。经验研究仍评估绝对损失契约，并未证明平方损失版本具有经验效用。推论~\ref{cor:timevarying_zh} 只在下端点条件成立时覆盖正的时变方向权重和非负预算。阈值型、非连续或其他非度量决策损失仍须重新推导精确可行集，本文也没有估计实际部署代价权重。第三，严格非增用于估计长期趋势，可能抑制真实的短时容量恢复或测量回弹，因此应同时保存原始观测与趋势输出。

第四，全部 $12$ 个开发域均已打开。嵌套设计避免当前外层目标标签选择 $\alpha$，预算路径协议也在该路径计算前完成锁定，但两者仍属于回顾性开发证据。第五，BaSyTec 流程在访问结果前经历两次有记录的结构重启：一个解析开发电芯被排除，且最终冻结前对 $47$ 个确认文件执行了只读表头普查。工作区能够证明预测先于容量访问，却不能证明团队范围从未历史接触结果；其 $45$ 电芯结论因此属于结果隔离的公开数据确认，而不是前瞻性实验室试验或独立实验室重复。第六，冻结参照范围排除了 $0\,^{\circ}\mathrm{C}$、$1.5C$ 工况的两个电芯；全工况敏感性虽保持结论，但属于事后分析。第七，NASA 压缩包只代表一个包含异构运行条件的公开实验室项目，只支持相对压力测试，不支持高精度绝对 SOH 或跨实验室普适性。第八，方法假定前五次充电记录能够构成有意义的投运参照；若电芯进入系统前已经发生未观测老化，或初始记录受到污染，该假设可能失效。

最后，PCHP 总会返回一个有界预测，并不判断何时应当拒绝输出。未来可以增加支持集或可估计性门控，但这将形成独立方法主张，必须在不使用结果的条件下预先定义，并使用一个在团队及项目谱系中结果均不可得的新验证面重新检验。

\section{结论}

本文针对目标结果与未来记录在预测时均不可得的跨域电池 SOH 场景，提出前缀因果损害预算投影。该方法把受保护源域预测转换为单向因果状态，再将投运参照学习候选投影到递归可行的损害管内，从而获得前缀不变性、非增输出、每条记录 $O(1)$ 的在线状态与计算，以及精确逐记录绝对损失损害界。同一分析还证明，该损害管是给定预算下必要且充分的最大输出集合，非平凡的无条件零损害更新不可能存在，并且 PCHP 在完整在线约束下唯一最小化候选失真。损失层分析进一步给出可复用的可行性边界：度量损失产生精确度量球；有界平方损失产生依赖基线位置且可沿用同一递归实现的精确区间；结果遍历 $\R$ 的平方损失则不允许任何非平凡有限结果一致预算。方向加权与时变推论进一步给出非对称损害管、下端点递归可行性的必要充分条件与紧的当前最小预算。覆盖 $12$ 个数据域的严格嵌套评估得到 $0.00486$ 的域等权 MAE 降幅；匹配的候选无关控制与协议锁定的回顾性预算路径共同表明，开发增益并非仅由某一个常数安全偏移或单一预算值解释。更强的无保护比较器取得显著更低的 MAE，却违反每个数据域的预算，从而明确了机会成本。结果隔离的 BaSyTec 确认在全部 $45$ 个合格电芯上改善并保持全部损害检查，而固定平移诊断说明该外部域验证的是契约，不是自适应优越性。NASA 压力测试则暴露绝对目标不兼容区域。因此，本文贡献是面向学习型电池健康更新的可部署信息与损害控制契约，而不是精度支配或普适 SOH 安全主张。

\section*{数据与代码可用性}

本文使用的精确开发数据快照 BatteryLife\_Processed v11 已由 Zenodo 以 CC BY~4.0 许可归档 \citep{batterylife2026processed}。其中各组成数据域仍遵循原始提供方各自的访问、署名与再分发条款，因此本文不随复现材料再分发任何第三方原始数据压缩包。包含 $48$ 个电芯的拜罗伊特退化路径数据集由 Zenodo 以 CC BY~4.0 发布 \citep{roder2025degradationdata}；NASA PCoE Battery Data Set~5 由 NASA Prognostics Data Repository 官方发布 \citep{saha2007nasa}。随投稿材料提供的可执行审稿轻量复现包包含完整代码闭包、获取说明、哈希、环境、预冻结记录、按域和电芯聚合的证据、确定性算子与统计重放检查及全部论文图生成脚本。该复现层不再分发第三方原始压缩包；完整重训与外部重放仍须按提供方条款自行下载数据。作者自有代码的发布许可证将在机构投稿信息确认时一并确定。

\section*{作者信息}

作者：吴雨阳。作者单位、通讯地址、ORCID、基金、致谢及利益冲突声明暂留空，待投稿信息最终确认后填写。

\bibliographystyle{plainnat}
\bibliography{references}

\end{document}
